\section{从武昌到退位：帝国在不同地方结束}

武昌起义使武汉三镇首先成为一处必须立刻处理兵员、枪械、渡江交通和城内秩序的战场。新军中的起事者和革命团体取得了部分城市空间，却还要面对清军的反攻、粮饷的筹集和消息向外传递的速度。朝廷所能看见的是一串战报、请饷和调兵的文书；城中与周边居民所面对的，则是军队驻扎、市场是否开张、家人能否通行以及哪一种命令会在次日仍然有效。这场起事因而不是一个抽象的“革命爆发”，而是把原先已因新军组织、革命网络和铁路风潮而绷紧的地方关系推入公开的军事竞争。

消息沿电报线、报纸和人际网络离开武汉，各省的回应却没有依同一张时刻表发生。有的省城先出现军队与官员的重新结盟，有的地方由商绅、咨议局人士和新式团体参与商议名义，有的地区则在省城宣布之后仍须逐县处理厘金、田赋、仓储和团练。所谓“独立”首先是一项政治宣言；它能否变成行政事实，要看谁能守住省城和交通线，谁能发给士兵军饷，谁能使税款不在征解途中断裂。旧督抚系统留下的印信、档册和征收网络可以被接管，也可以被地方军队、县衙与商会各自截留。因而，省份脱离北京的次序不能直接当作地方社会已经换政权的次序。

清廷起用袁世凯处理北方军政和议和事务，是在中央已难以单凭宫廷命令调动可信军队的条件下作出的选择。北洋军的指挥、北京的安全和与南方交涉的门径随之被放到同一张谈判桌上，但它们并不天然指向同一个结果。袁世凯既须向清廷说明军政处置，又要同革命方面周旋；革命方面也不能只以宣言取代军队、港口和城市财政的安排。战事与议和互相塑形：战场上的推进会改变谈判筹码，谈判的风声也会影响官员、士兵和地方团体究竟继续观望还是转向。

南京临时政府成立时，南北谈判尚未结束。它提出新的法理中心，却不能因此抹平各省军政和财政的实际差异。铁路国有化与外债引起的冲突已经显示，铁路股权、借款责任和地方筹款并非单纯的技术问题；辛亥之际，盐税、关税、厘金以及地方所收各项税款由谁保管、用来支付何处军队、是否承认旧债，同样决定了新名义能在何处落地。省级机关的改名、旧官的去留和军队的改编常在同一座城市里并行发生，不能把一纸任命或一份公报当成已经完成的交接。

北部边疆的时间也不服从北京与南京之间的叙事节奏。库伦方面在1911年末宣布脱离清廷，并拥立哲布尊丹巴为政教领袖；这一行动发生在清廷的财政、军政和通信同时承压之际，也牵动蒙古王公、宗教网络和跨境关系。公告能够表明一种政治主张，却不能单独画出稳定的控制范围，更不能把各地的支持写成整齐一致。青藏高原、新疆以及东北、西南的政局、交通与地方权威，各有尚待按本地材料厘清的转换过程。帝国的结束在边地不是从一份中央文件向外铺开的直线，而是多种旧关系和新主张在不同道路、城镇与文书语言中重新碰撞。

1912年的退位诏书为皇帝、部院和旗务的旧名义划出了明确的制度终点。南北谈判、财政军政压力和皇室待遇的安排共同促成这一结果；诏书所写的优待可以确认中央层面的承诺，却不能替代对各地兑现情形的追查。战争、债务、地方军队和身份供养并未随诏书消失。旗人家庭要面对生计与旧有供养的变动，旧官吏要决定是否留任，地方学校、商会和新式军队也要在新的制度名义下重新寻找资源与位置。对于许多县乡居民而言，最先可见的变化未必是国家称号，而可能是征粮者、驻军者、县署告示和市场秩序的改变；这些变化并不在所有地方同时到来，也不必然朝同一方向发展。

所以，帝国终结不是末代皇帝离开权位的一瞬，而是一组行政、军事、财政和社会关系被拆开重组的开端。退位解决了谁在中央名义上继承国家的问题，却把军队的归属、税源的流向、地方债务与边地权威留给此后的共和国继续承受。

\subsection{材料与争议}

退位诏书、内阁档案、各省公报与临时政府文书能把程序、命令和自我表述放回日序，却只能首先证明这些文书被拟定、传递或刊布。革命宣传、私人日记、报刊与外交报告保存的是动员者或观察者各自所处的位置；它们有助于看见消息如何流动，不能自动量出一省的支持程度。地方档案、商会账簿、税册和军队给饷记录更接近日常交接，但也常只留下一段机构的视野。

因此，省份独立的日期、地方实际控制、军队人数和群众态度须逐项互校。尤其是关于边疆，清方、俄方和蒙古文材料使用的称谓与历法并不一致；关于城市与县乡，电报抵达、报刊刊登、官署改名和税粮实际改征也不是同一件事。材料所能支持的是不同地点的局部时序，而不是一幅全国同步转向的图景。

\subsection{交叉阅读窗口：少年中国的排比，和1911年的不同时钟}

梁启超1900年刊行的《少年中国说》写道：“少年智则国智，少年富则国富。”这是一篇流亡改革派的政治散文，整齐排比把国家的将来托付给“少年”，以夸张、召唤式的节奏制造共同体的期待；它并非1911年武昌、四川或北京的现场报道，更不能替所有学生、士兵和乡村居民规定政治意愿。把它同《宣统政纪》、内阁档、各省公报和临时政府文书并读，能钉住政策、退位与交接程序；四川保路档案、商会账簿、地方报刊、《申报》与《民立报》则显示消息、产权和动员如何在不同城市流转，英日使领馆报告提供另一种受外交利益限制的观察。

钱穆的提问在于立宪、内阁和省级机构为何未能把新旧权力的交接稳定为共同程序；吕思勉的提问则在于债务、盐税、铁路股权、教育与地方组织怎样决定哪些人有条件进入公共政治。报刊可证明某日消息被刊登和以何种语气出现，不能证明全省同时知晓；电报和独立宣言也不等于县乡已经易手。梁启超的散文让我们听到革命前公共语言的音量，而多源互校才允许我们辨认帝国是在何处、以何种行政和社会节奏结束。
